% Standalone problem document, generated from the adjacent README.md.
% Compile with XeLaTeX. Mathematical content is maintained in Markdown.
\documentclass[11pt,a4paper]{article}
\usepackage[fontsize=10.5pt]{scrextend}
\usepackage[margin=22mm,headheight=15pt,headsep=9mm,footskip=11mm]{geometry}
\usepackage{fontspec}
\setmainfont{texgyrepagella}[Extension=.otf,UprightFont=*-regular,BoldFont=*-bold,ItalicFont=*-italic,BoldItalicFont=*-bolditalic]
\setsansfont{texgyreheros}[Extension=.otf,UprightFont=*-regular,BoldFont=*-bold,ItalicFont=*-italic,BoldItalicFont=*-bolditalic]
\setmonofont{texgyrecursor}[Extension=.otf,UprightFont=*-regular,BoldFont=*-bold,ItalicFont=*-italic,BoldItalicFont=*-bolditalic,Scale=MatchLowercase]
\usepackage{amsmath,amssymb}
\usepackage{unicode-math}
\setmathfont{texgyrepagella-math.otf}
\usepackage{xcolor,microtype,enumitem,needspace,fancyhdr}
\usepackage{bookmark}
\definecolor{ink}{HTML}{17344B}
\definecolor{accent}{HTML}{197D80}
\definecolor{muted}{HTML}{596773}
\hypersetup{colorlinks=true,linkcolor=accent,urlcolor=accent,pdftitle={MF-14 - Degree
coverage with seven matrix
multiplications},pdfauthor={Open Problems in Numerical Linear Algebra}}
\urlstyle{same}
\setlength{\parindent}{0pt}
\setlength{\parskip}{5pt plus 1pt minus 1pt}
\linespread{1.04}
\setlength{\emergencystretch}{3em}
\widowpenalty=10000
\clubpenalty=10000
\displaywidowpenalty=10000
\allowdisplaybreaks[1]
\setlist{nosep,leftmargin=1.5em}
\providecommand{\tightlist}{\setlength{\itemsep}{0pt}\setlength{\parskip}{0pt}}
\providecommand{\pandocbounded}[1]{#1}
\pagestyle{fancy}
\fancyhf{}
\fancyhead[L]{\sffamily\footnotesize\color{muted}OPEN PROBLEMS IN NUMERICAL LINEAR ALGEBRA}
\fancyhead[R]{\sffamily\bfseries\footnotesize\color{accent}MF-14}
\fancyfoot[L]{\parbox[b]{0.9\textwidth}{\sffamily\fontsize{8}{10}\selectfont\color{muted}Editorial ratings; a bounded search cannot certify that a problem remains unsolved.\\Literature check: 2026-09-19}}
\fancyfoot[R]{\sffamily\footnotesize\color{muted}\thepage}
\renewcommand{\headrulewidth}{0.3pt}
\renewcommand{\headrule}{\hbox to\headwidth{\color{accent}\leaders\hrule height \headrulewidth\hfill}}
\makeatletter
\renewcommand{\section}{\Needspace{5\baselineskip}\@startsection{section}{1}{0pt}{12pt plus 2pt minus 2pt}{4pt}{\sffamily\bfseries\large\color{ink}}}
\renewcommand{\subsection}{\Needspace{4\baselineskip}\@startsection{subsection}{2}{0pt}{9pt plus 1pt minus 1pt}{3pt}{\sffamily\bfseries\normalsize\color{ink}}}
\makeatother
\setcounter{secnumdepth}{0}
\begin{document}
{\sffamily\small\color{muted}Matrix functions and stability\par}
\vspace{3pt}
{\raggedright\hyphenpenalty=10000\exhyphenpenalty=10000\sffamily\bfseries\fontsize{21}{24}\selectfont\color{ink}Degree
coverage with seven matrix multiplications\par}
\vspace{7pt}
{\sffamily\small\textbf{Difficulty:} challenging\quad\textbf{Importance:} interesting
to the community\par}
{\sffamily\small\bfseries\color{accent}Status: Lean verified\par}
\vspace{4pt}
{\color{accent}\hrule height 0.8pt}
\vspace{6pt}
\textbf{Rating rationale:} Historical ratings for the original
conjecture: challenging because numerical evidence for coefficient-space
coverage needs exact algebraic justification; community impact comes
from multiplication budgets for practical matrix-function evaluation.

\subsection{Lean proof and verification
evidence}\label{lean-proof-and-verification-evidence}

\textbf{Lean verified --- 19 September 2026 (UTC). Formalization: George
Stepaniants}, Department of Computing and Mathematical Sciences,
California Institute of Technology, with substantial AI assistance.
Marcus Webb, The University of Manchester retains mathematical
construction authorship.

Seven chronological products cover every complex polynomial through
degree 44 in the full 129-coefficient closure model, refuting the
original equality 42. The later exact maximum 47 is outside this
formalization.

All 25 frozen contracts passed actual local serial Lean, two wholly
nonauthor final AI-agent source reviews, and the
\href{https://github.com/sgstepaniants/OpenProblemsInNLA/actions/runs/35433896875}{sgstepaniants
run 35433896875} and
\href{https://github.com/ajt60gaibb/OpenProblemsInNLA/actions/runs/35433898434}{ajt60gaibb
run 35433898434}. The real Comparator, default kernel, standard-axiom
checks, sandbox and rejection controls passed on the immutable proof
\href{https://github.com/sgstepaniants/OpenProblemsInNLA/tree/c89c859a6b21c399a99f5e9f5fb2b7f472c96657/matrix-functions-and-stability/MF-14/lean}{c89c859a}.
Two independent statement reviews preceded implementation. Kernel-mode
LeanCert audits are retained, with exact computation minimized. An
independent operational audit authenticates the execution evidence. No
official Tau Ceti service or external human peer review is claimed.

\href{https://github.com/ajt60gaibb/OpenProblemsInNLA/blob/main/matrix-functions-and-stability/MF-14/lean/README.md}{Proof,
reproduction and scope} ·
\href{https://github.com/ajt60gaibb/OpenProblemsInNLA/blob/main/matrix-functions-and-stability/MF-14/lean/verification/linux-2026-09-19/SUMMARY.json}{Exact
execution evidence} ·
\href{https://github.com/ajt60gaibb/OpenProblemsInNLA/blob/main/matrix-functions-and-stability/MF-14/lean/formalization.yaml}{Formalization
metadata} ·
\href{https://github.com/ajt60gaibb/OpenProblemsInNLA/pull/305}{Upstream
PR \#305}.

\subsection{Resolution --- exact maximum 47, 17 September
2026}\label{resolution-exact-maximum-47-17-september-2026}

\textbf{Author:} Marcus Webb, The University of Manchester. Developed
with substantial ChatGPT/Codex assistance. Separate Codex agents
reviewed the mathematics, and an independently implemented exact
verifier reconstructed the full polynomial and its derivatives. This is
informal AI-agent review, not external human peer review or formal
verification.

\href{https://github.com/ajt60gaibb/OpenProblemsInNLA/blob/main/references/webb-mf14-degree47-2026-09-17/proof.pdf}{Theorem
1 of the complete proof} determines the exact maximum in the retained
question:

\[d_7=47.\]

Every complex polynomial of degree at most 47 belongs to the Zariski
closure of seven-product outputs in the full coefficient space through
degree 128. An explicit seven-product family has a rigorously certified
order-47 contact with an invertible 48-by-48 coefficient Jacobian. Exact
Gaussian-rational contraction bounds prove the contact exists; the
inverse function theorem and a weighted input/output limit give the
universal coverage statement, with every higher coefficient tending to
zero.

A self-contained available-space dimension count gives
\(\dim\overline{\mathcal P_7}^{\,Z}\le49\). Irreducibility and the
seven-product polynomial \(x^{128}\) then exclude coverage of the entire
degree-at-most-48 space. The new proof does not depend on the earlier
border construction or on an external dimension theorem. Exact
representation of every degree-47 polynomial and numerically stable
coefficient recovery are not asserted.

\href{https://github.com/ajt60gaibb/OpenProblemsInNLA/blob/main/references/webb-mf14-degree47-2026-09-17/proof.tex}{Standalone
proof source} ·
\href{https://github.com/ajt60gaibb/OpenProblemsInNLA/blob/main/references/webb-mf14-degree47-2026-09-17/README.md}{Exact
certificates and reproduction} ·
\href{https://github.com/ajt60gaibb/OpenProblemsInNLA/blob/main/references/webb-mf14-degree47-2026-09-17/verification/packaging-review.md}{Independent
review}. The original equality \(d_7=42\), permanent ID and canonical
path remain unchanged. The earlier degree-42 and degree-44 results below
retain their attribution and archived evidence.

\subsection{Earlier negative resolution --- degree 44, 17 September
2026}\label{earlier-negative-resolution-degree-44-17-september-2026}

\textbf{Author:} Marcus Webb, The University of Manchester. Developed
with substantial ChatGPT assistance; two separate Codex agents
independently reviewed the proof and recomputed its exact certificates.
This is informal AI-agent review, not external human peer review or
formal verification.

The conjectured equality is false.
\href{https://github.com/ajt60gaibb/OpenProblemsInNLA/blob/main/references/webb-mf14-degree44-2026-09-17/proof.pdf}{Theorem
1 of the complete proof} establishes

\[\mathbb C[x]_{\le44}\subseteq\overline{\mathcal P_7}^{\,Z}.\]

The construction first retains a suitable quadruple of polynomials in
the joint closure of four-product computations (Lemma 2). Three more
products give a 45-parameter coefficient map whose exact Jacobian
determinant is 256. The argument uses the original complex field and all
coefficients through degree 128; high-degree coefficients vanish in the
limit rather than being discarded.

Corollary 3, using Jarlebring--Lorentzon's dimension theorem, gives
\(44\le d_7\le47\) for the maximum in the retained question. That
contribution left the exact maximum undetermined while giving a complete
negative answer to the specific equality \(d_7=42\); the new theorem
above now determines it as 47. No exact representation of every
degree-44 polynomial, real Euclidean density, or numerical stability is
asserted.

\href{https://github.com/ajt60gaibb/OpenProblemsInNLA/blob/main/references/webb-mf14-degree44-2026-09-17/proof.tex}{Proof
source} ·
\href{https://github.com/ajt60gaibb/OpenProblemsInNLA/blob/main/references/webb-mf14-degree44-2026-09-17/README.md}{Independent
reviews and reproduction} ·
\href{https://github.com/ajt60gaibb/OpenProblemsInNLA/blob/main/references/webb-mf14-degree44-2026-09-17/README.md\#source-and-attribution}{Source
conversation and attribution}. The permanent ID, canonical path and
original mathematical target are retained. The degree-42 result below
retains its separate credit.

\subsection{Verified partial result ---
2026-09-11}\label{verified-partial-result-2026-09-11}

\textbf{Author:} Matthew J. Colbrook, Department of Applied Mathematics
and Theoretical Physics, University of Cambridge. \textbf{Independent
Codex-agent review: PASS for the stated partial results.}

A fixed seven-product scheme has a full-rank complex coefficient map,
certified by a nonzero exact integer Jacobian minor. Its image contains
a nonempty Zariski-open subset of \(\mathbb C[x]_{\le42}\) and is
Euclidean dense there.

This earlier result did not prove an upper bound excluding degree 43 and
above, so at the time it left the maximal-degree equality open. It
asserted neither exact representation of every polynomial nor real
Euclidean dense coverage. The later degree-44 construction above now
refutes that equality; the degree-42 theorem remains valid.

\textbf{Primary reference:}
\href{https://github.com/ajt60gaibb/OpenProblemsInNLA/blob/main/references/colbrook-matrix-functions-2026-09-11/manuscripts/MF-14.pdf}{complete
authored PDF},
\href{https://github.com/ajt60gaibb/OpenProblemsInNLA/blob/main/references/colbrook-matrix-functions-2026-09-11/manuscripts/MF-14.tex}{standalone
TeX}, \textbf{Theorem 1 and equations (1)-(2)}.
\href{https://github.com/ajt60gaibb/OpenProblemsInNLA/blob/main/references/colbrook-matrix-functions-2026-09-11/verification/reviews/MF-14-review.md}{Independent
proof review} ·
\href{https://github.com/ajt60gaibb/OpenProblemsInNLA/blob/main/references/colbrook-matrix-functions-2026-09-11/README.md}{Authorship
and submission record}. Verification is independent agent review, not
external human peer review or formal certification.

\newpage

\subsection{Problem statement}\label{problem-statement}

Start with the scalar polynomials \(1,x\in\mathbb C[x]\). Allow
arbitrarily many complex linear combinations of already computed
polynomials and at most seven multiplications of two already computed
polynomials. Let \(\mathcal P_7\subseteq\mathbb C[x]_{\le128}\) be the
set of possible outputs.

Identify a polynomial with its coefficient vector in \(\mathbb C^{129}\)
and let \(\overline{\mathcal P_7}^{\,Z}\) be its Zariski closure: the
common zero set of all polynomial equations in the coefficients that
vanish throughout \(\mathcal P_7\). Is

\[\max\{d\in\{0,\ldots,128\}:\mathbb C[x]_{\le d}\subseteq\overline{\mathcal P_7}^{\,Z}\}=42?\]

The target concerns closure and the complex coefficient field. Exact
representability of every polynomial is a stronger assertion.

\subsection{Why it matters}\label{why-it-matters}

For a dense input matrix, matrix products dominate the cost of
polynomial evaluation. The conjecture determines how much polynomial
degree a budget of seven products can cover, providing a theoretical
benchmark for matrix-function evaluation schemes.

\subsection{References}\label{references}

\begin{enumerate}
\def\labelenumi{\arabic{enumi}.}
\tightlist
\item
  E. Jarlebring and G. Lorentzon, \emph{The polynomial set associated
  with a fixed number of matrix-matrix multiplications},
  arXiv:2504.01500v3 (13 August 2025), §2.2 (closure convention), §5.4,
  Conjecture 13. \href{https://arxiv.org/html/2504.01500}{Primary text}.
\item
  J. Sastre, J. Ibáñez, J. M. Alonso, and E. Defez, \emph{Beyond
  Paterson--Stockmeyer: Advancing matrix polynomial computation}, WSEAS
  Transactions on Mathematics 24 (2025), 684--693, §4 (a five-product
  construction).
  \href{https://wseas.com/journals/mathematics/2025/b385106-036\%282025\%29.pdf}{Primary
  paper}, \href{https://doi.org/10.37394/23206.2025.24.68}{DOI}.
\item
  J. M. Alonso, J. Sastre, J. Ibáñez, and E. Defez, \emph{A systematic
  framework for stable and cost-efficient matrix polynomial evaluation},
  arXiv:2603.23143v1 (24 March 2026), abstract and conclusions.
  \href{https://arxiv.org/html/2603.23143}{Primary text}.
\end{enumerate}

\subsection{Status check --- 2026-09-08}\label{status-check-2026-09-08}

The latest arXiv source remains v3 and labels the displayed equality a
conjecture supported by computations. The later papers give
constructions and stability procedures without establishing the
seven-product degree-coverage claim. Searches included the exact title
with \textit{2026},
\textit{Jarlebring Lorentzon conjecture seven multiplications 42},
and \textit{Sastre polynomial evaluation 2026}. No resolution was
located. A
\href{https://indico3.mpi-magdeburg.mpg.de/event/58/contributions/1074/}{3
September 2026 author talk} also describes higher-cost degree coverage
as conjectural. Six-product real/complex formulations have a separate
source ambiguity and are deliberately not included here.

\subsection{Audit --- 2026-09-10}\label{audit-2026-09-10}

Rechecked \href{https://arxiv.org/html/2504.01500}{Conjecture 13 and
Remark 14}: complex closure coverage at seven products is still
conjectural, with concrete matrix-exponential applications. Title and
seven-product searches found no exact proof. The
\href{https://arxiv.org/html/2606.24701v1}{June 2026 degree-eight
construction} solves a different budget/degree problem.
\end{document}
